\section*{Chapitre \ref{ch:dv:local-volatility} --- Volatilité locale}

\begin{solution}{exo:dv:local-volatility:1}
$\sigma_{\mathrm{loc}}=\sqrt{0.045/0.9}=\sqrt{0.05}=22.4\%$.
\end{solution}

\begin{solution}{exo:dv:local-volatility:2}
$16.4\%+2\times(-0.47)\times(-0.1)=25.7\%$, contre 27.4~\%~: la règle capture la pente~; le reste est la
courbure du smile, que la \omterm{def:dv:local-volatility:model}{volatilité locale} double également.
\end{solution}

\begin{solution}{exo:dv:local-volatility:3}
Indépendante du spot~: $\sigma_{\mathrm{loc}}^2=0.5\times0.15^2+0.5\times0.30^2=0.05625$, soit une volatilité plate de 23.7~\%. En 90
avec une probabilité de 0.8 de l'état haut~: $0.8\times0.09+0.2\times0.0225=0.0765$, 27.7~\%.
\end{solution}

\begin{solution}{exo:dv:local-volatility:4}
La formule divise le calendar spread par le butterfly~: un \omterm{def:dv:implied-volatility-and-its-surface:static}{arbitrage calendaire} rend la variance locale
négative, un \omterm{def:dv:implied-volatility-and-its-surface:static}{arbitrage butterfly} rend le dénominateur négatif ou nul, et l'un ou l'autre ne donne aucune \omterm{def:dv:local-volatility:model}{volatilité locale}.
Une cotation bruitée déforme localement le smile, et la dérivée seconde au dénominateur amplifie cette déformation en une
pic ou trou dans la \omterm{def:dv:local-volatility:model}{volatilité locale} au voisinage de la cotation.
\end{solution}

\begin{solution}{exo:dv:local-volatility:5}
$\sqrt{(0.1919^2\times1-0.1779^2\times0.5)/0.5}=20.49\%$, contre 20.68~\% pour l'option à départ différé à la
monnaie du modèle~: le niveau du \omterm{def:dv:local-volatility:forward}{smile forward} est essentiellement celui de la variance forward~; le petit écart tient à la courbure du
\omterm{def:dv:local-volatility:forward}{smile forward} à la monnaie.
\end{solution}

\begin{solution}{exo:dv:local-volatility:6}
$2\times0.47\times0.01=0.94$ point, contre 1.31 simulé~: la règle tient pour les échéances courtes et les skews faibles, et
ce skew à trois mois est pentu.
\end{solution}

\begin{solution}{exo:dv:local-volatility:7}
Avec 50\,000 trajectoires, l'erreur du put 90 est de $-0.04$ point, avec une erreur type de 0.13~; avec 200\,000, $+0.11$, avec une
erreur type de 0.06. Les erreurs de tous les strikes à 200\,000 trajectoires ont le même signe~: un biais de discrétisation du
pas de temps quotidien (environ 0.08 point), biais que davantage de trajectoires ne fait pas disparaître mais que quatre pas par jour éliminent.
\end{solution}

\begin{solution}{exo:dv:local-volatility:8}
Un paiement européen sur l'indice à une date, oui~: son prix ne dépend que de la distribution marginale à cette date,
que la calibration reproduit. Une option à départ différé paie sur le ratio de l'indice à deux dates, ce qui dépend de la
distribution conjointe, donc de la dynamique du modèle~: la \omterm{def:dv:local-volatility:model}{volatilité locale} lui donne un \omterm{def:dv:local-volatility:forward}{smile forward} plus plat que
celui d'un modèle à volatilité stochastique calibré sur les mêmes options vanilles.
\end{solution}

\begin{solution}{pb:dv:local-volatility:1}
\textbf{1.} 16.4~\% et 17.7~\% à la monnaie~; 21.1~\% et 27.4~\% à $k=-0.1$.
\textbf{2.} 1.98.
\textbf{3.} La volatilité implicite est la moyenne de la \omterm{def:dv:local-volatility:model}{volatilité locale} entre le spot et le strike~; pour produire une
pente donnée de cette moyenne, la \omterm{def:dv:local-volatility:model}{volatilité locale} doit avoir une pente deux fois plus forte.
\textbf{4.} $+0.08$, $+0.06$, $+0.09$ point de volatilité.
\textbf{5.} Un problème de grille ou d'extrapolation, ou un arbitrage dans la surface en entrée, et non du bruit.
\textbf{6.} 20.7~\% et 17.8~\%.
\textbf{7.} La volatilité forward à la monnaie de six mois à un an vaut 20.5~\% d'après la structure par terme~: le
\omterm{def:dv:local-volatility:forward}{smile forward} se situe au niveau de la variance forward.
\textbf{8.} $-0.19$ et $-0.34$ par unité de \omterm{def:dv:implied-volatility-and-its-surface:surface}{log-moneyness}.
\textbf{9.} 0.57.
\textbf{10.} Sa \omterm{def:dv:local-volatility:model}{volatilité locale} est forte près du spot du jour aux temps courts~; six mois plus tard, le spot s'est dispersé et
les \omterm{def:dv:local-volatility:model}{volatilités locales} pertinentes sont celles des temps plus longs, où la fonction est plus plate.
\textbf{11.} 21.8~\%~; la forme actuelle, décalée au niveau forward, donne 22.5~\%.
\textbf{12.} 3.82 et 4.00 pour 100.
\textbf{13.} La \omterm{def:dv:local-volatility:model}{volatilité locale}, le modèle le moins cher pour la protection qu'elle vend~; un modèle qui valorise le risque plus bas
est un risque de modèle (chapitre 27).
\textbf{14.} Des options à départ différé ou des cliquets cotés par les dealers, ou la variance des skews implicites futurs observés
historiquement, qui révèlent le comportement des skews.
\textbf{15.} Valoriser avec les deux modèles et conserver l'écart, 0.18 pour 100, en réserve de modèle jusqu'à ce que le risque disparaisse.
\textbf{16.} 23.4~\% sous \omterm{def:dv:local-volatility:model}{volatilité locale}, 18.8~\% sous \omterm{def:dv:implied-volatility-and-its-surface:sticky}{sticky strike}, à partir de 16.4~\%.
\textbf{17.} Pour les produits dont la valeur dépend des distributions marginales et de la dynamique aux échéances courtes (barrières proches
du spot, options américaines), et comme la \omterm{def:dv:local-volatility:projection}{projection markovienne} à laquelle les autres modèles sont calibrés.
\textbf{18.} Un aléa sur la volatilité qui n'est pas fonction du spot~: volatilité stochastique (chapitre 10) ou un
modèle de variance forward (chapitre 12).
\textbf{19.} 0.57~: le skew à six mois, six mois plus tard, vaut 57~\% du skew à six mois du jour.
\textbf{20.} La dynamique~: comment le smile se déplace avec le spot et à quoi il ressemblera aux dates futures.
\end{solution}

\begin{solution}{iq:dv:local-volatility:1}
La volatilité, fonction déterministe du temps et du spot, $\sigma_{\mathrm{loc}}(t,S)$, d'une diffusion qui
reproduit le prix de toutes les options européennes. Elle est unique parce que les prix des calls déterminent la densité à chaque date
(Breeden--Litzenberger), et l'évolution de la densité détermine le coefficient de diffusion (\omterm{thm:dv:local-volatility:dupire}{formule de Dupire}).
\iqlookfor{des prix aux densités, puis au coefficient.}
\end{solution}

\begin{solution}{iq:dv:local-volatility:2}
Elle monte, d'environ deux fois ce que donne le \omterm{def:dv:implied-volatility-and-its-surface:sticky}{sticky strike}~: la volatilité implicite à la monnaie est proche de la
\omterm{def:dv:local-volatility:model}{volatilité locale} au spot, qui a une pente deux fois plus forte que le skew implicite, et le spot s'est déplacé vers une zone de
\omterm{def:dv:local-volatility:model}{volatilité locale} plus élevée.
\iqlookfor{le facteur deux et sa raison.}
\end{solution}

\begin{solution}{iq:dv:local-volatility:3}
La densité $q$ satisfait $\partial_Tq=\tfrac12\partial_{ss}(\sigma^2s^2q)$~; $C=\int(s-K)^+q\,ds$~; après deux intégrations par parties, $\partial_TC=
\tfrac12\int(s-K)^+\partial_{ss}(\sigma^2s^2q)\,ds=\tfrac12\sigma^2(T,K)K^2q(T,K)$~;
$q=\partial_{KK}C$. D'où $\sigma^2(T,K)=2\partial_TC/(K^2\partial_{KK}C)$.
\iqlookfor{l'équation de Kolmogorov progressive et les intégrations par parties.}
\end{solution}

\begin{solution}{iq:dv:local-volatility:4}
Ils coïncident sur tout ce qui ne dépend que de la distribution du sous-jacent à des dates isolées (paiements européens,
y compris les digitales et les swaps de variance sans sauts). Ils peuvent différer sur tout ce qui dépend de la loi conjointe à
plusieurs dates ou des trajectoires~: barrières, asiatiques, lookback et options à départ différé, cliquets, autocallables.
\iqlookfor{marginales contre lois conjointes.}
\end{solution}

\begin{solution}{iq:dv:local-volatility:5}
Ces instabilités viennent de la dérivation d'une surface bruitée ou sujette à arbitrage, des cassures de pente lors de l'interpolation entre les tranches, et de
l'extrapolation au-delà des strikes cotés. Calibrer d'abord une surface paramétrique lisse sans arbitrage, calculer la
\omterm{def:dv:local-volatility:model}{volatilité locale} à partir de celle-ci en variance totale, la borner par un cap et un floor là où l'entrée ne contient aucune information, et vérifier par revalorisation.
\iqlookfor{lisser l'entrée, non la sortie.}
\end{solution}

\begin{solution}{iq:dv:local-volatility:6}
Le \omterm{def:dv:local-volatility:model}{modèle de volatilité locale} avec $\sigma^2(t,K)=\E[\sigma_t^2\mid S_t=K]$, qui a les mêmes lois marginales que le
modèle à volatilité stochastique. Pour calibrer un modèle à volatilité stochastique sur toute la surface, multiplier sa volatilité par un
fonction de levier $L(t,S)$ choisie de sorte que $L^2(t,K)\,\E[\sigma_t^2\mid S_t=K]$ égale la variance locale du marché~: la
modèle à \omterm{def:dv:local-volatility:model}{volatilité locale} stochastique du chapitre 20.
\iqlookfor{espérance conditionnelle de la variance sachant le spot, et fonction de levier.}
\end{solution}
